\begin{table}[H]
\centering
\caption[Component ablation settings]{Component ablations analyzed against Final-HV reward. Each ablation preserves the Final-HV quality signal and changes one scheduler or reward-side component.}
\label{tab:ablation-settings}
\small
\renewcommand{\arraystretch}{1.12}
\setlength{\tabcolsep}{5pt}
\begin{tabular}{>{\raggedright\arraybackslash}p{0.25\linewidth}>{\raggedright\arraybackslash}p{0.31\linewidth}>{\raggedright\arraybackslash}p{0.34\linewidth}}
\toprule
Ablation & Disabled component & Hypothesis tested \\
\midrule
No budget annealing & The remaining-budget decay factor in the cluster-bandit exploration term is disabled, so exploration pressure no longer decreases as the budget is consumed. & Tests whether annealing exploration toward the end of a run improves budget efficiency and prevents late-budget over-exploration. \\
No Page-Hinkley & The Page--Hinkley drift threshold is set high enough that cluster-arm reset is effectively disabled in the analyzed runs. & Tests whether resetting stale cluster statistics helps when cluster reward distributions change during heuristic-population search. \\
No diversity reward & The positive cumulative-diversity-gain term is removed from the scalar reward by setting the diversity weight to zero. Population management, non-dominated sorting, and crowding-distance truncation remain unchanged. & Tests whether explicitly rewarding expansion of the heuristic score distribution improves exploration, budget efficiency, or terminal population quality. \\
\bottomrule
\end{tabular}
\end{table}
